\section{Introduction}\label{sec:intro}

\paragraph{The problem.} A \emph{storage code} for a recovery relation $N$ on a finite coordinate set $V$ is a set of
words $\mathcal C\subseteq\mathcal A^V$ in which every symbol $x_q$ is determined by the symbols on its recovery set
$N(q)$. Storage codes are the fibres of index codes, and their rates measure how much a common broadcast can be
compressed given side information. The series of papers to which this one belongs studies whether a deterministic
multitape Turing machine can transpose a $d\times d$ matrix of $b(d)=\max(1,\lceil\log_2d\rceil)$-bit entries in time
linear in $L=d^2b(d)$, the statement~(T). For oblivious machines, the fifth paper reduced a large part of the
impossibility question to a machine-free coding question~\cite{DirectedRecovery2026}. The actual recovery relations of
linear-time computations have bounded incidence degree after cheap trimming, and only $o(\abs V)$ coordinates have short
\emph{paired distance} $\dist(L_q,R_q)$ between the two copies of $q$ in the bipartite incidence graph. Question~9.1 of
that paper asks whether these two properties alone force every storage code to have rate at most $1/2+o(1)$
(\cref{prob:q91}). An affirmative answer would bound the rate of all finite-alphabet codes on the cone relations of oblivious
linear-time machines~\cite[Proposition~9.2]{DirectedRecovery2026}, and several later papers of the series state results
conditional on it~\cite{BoundedNesting2026,DirectedReciprocity2026,MergeDepth2026}.

\paragraph{The results.} The answer is negative.
\begin{itemize}
\item \textbf{Theorem~A} (\cref{thm:A}). For $p=127$ there are recovery relations of incidence degree at most
$\Delta<127^{8\,791\,448}$, on which every paired distance tends to infinity, carrying $\F_{127^2}$-linear storage codes of
rate at least $8100/16130>0.5021$.
\item \textbf{Corollary~A$'$} (\cref{cor:Aprime}). Writing each symbol in $14$ bits gives binary storage codes of rate
greater than $0.5013$ on recovery relations of incidence degree at most $14\Delta$ with the same property. So the fixed
binary alphabet, singled out in Question~9.1, has a negative answer too.
\item \textbf{Theorem~A$^*$} (\cref{thm:Astar}). For all large primes $p$ inert in $\Q(\sqrt5)$, the guaranteed rates come
arbitrarily close to $(\sqrt3-1)^2=4-2\sqrt3\approx0.536$.
\item \textbf{Corollary~U} (\cref{cor:U}). Forgetting the directions of the relations of Theorem~A gives graphs of maximum
degree at most $2\Delta$ in which every odd cycle is long, carrying the same codes, and binary codes of rate greater than
$0.5013$ on graphs of degree at most $28\Delta$. So the undirected form of the question, the fourth paper's undirected
crux~\cite[Question~6.1]{CurvedMotion2026}, has a negative answer too.
\end{itemize}

The relations are \emph{twisted quotients} of the Bruhat--Tits tree $\Tree$ of $\PGL_2(\Q_{p^2})$, and the codes are
Tanner codes~\cite{Tanner1981} whose local codes are Reed--Solomon codes of length $p^2+1$. Theorem~T (\cref{thm:T})
separates the argument into two halves.
\begin{itemize}
\item \emph{The combinatorial half.} \textbf{Theorem~D} (\cref{thm:D}) shows that every $\delta$-regular graph on $n$
vertices has a \emph{lift selection} (one lift of each edge in a double cover) of degeneracy
$(2-\sqrt3)\delta+O(\sqrt{\delta\log n})$, through an explicit balanced charging rule. On the tree, such a selection splits
the edges into two classes, exchanged by a character $\chi$ of the group. Every edge then has a \emph{private check}, a
finite sum of local parity checks that sees no other edge of its class; we construct it explicitly by \emph{peeling}
(\cref{lem:peeling}). The resulting recovery relation never connects the two copies of a coordinate on the tree. On a
quotient, a path between them must lift to a path from an edge to a translate in the other class, which is far away when
the quotient has large girth (\cref{prop:descent}).
\item \emph{The algebraic half.} \textbf{Theorem~G} (\cref{thm:G}) shows that the Tanner codes of an explicit twisted
datum over $\F_{p^2}$ on the congruence quotients $\Gamma^+(\lambda_0^j)\backslash\Tree$ of the icosian group have dimension at
least $(k-1)^2/(p^2+1)$ times their length, for every weight $2\le k\le p-2$. These codes are spaces of sections of the
bundle of parallel weight $k$ on the Goren--Oort divisor of a Hilbert modular surface in characteristic~$p$, whose dual
graph is the quotient of the tree~\cite{BachmatGoren1999,KudlaRapoport1999,TianXiao2016}; Riemann--Roch and the Hasse
invariant give their dimension. Counting constraints would give less than one half.
\end{itemize}
At $p=127$ and $k=91$ the two halves just fit: \cref{thm:D} needs local redundancy $t\ge4640$, the weight gives
$t=127^2-91\cdot126=4663$, and the rate is $(91-1)^2/(127^2+1)>1/2$.

\paragraph{How to read the paper and its formal proofs.} Every numbered result states its status. A line \textsf{Lean:}
names the Lean declarations of a formalized result and their file, and the result is followed by an outline of the formal
proof; the Lean statement is authoritative. A line \textsf{Lean (partial):} names what is formalized and what is not.
Results marked \textsf{Not formalized} have complete proofs in the text. Remarks labelled as heuristic or numerical are not used in
any proof.

The combinatorics is formalized in Lean~4 with Mathlib, on the standard axioms only: the deterministic part of
\cref{thm:D} (options, rounding by Hall's theorem, and the fluid identities), the peeling lemma, the validity of checks on
Tanner codes and the interpolation property of MDS local codes, the descent to quotients (storage, incidence degrees,
paired distances), the binary blow-ups, the symmetrized graphs of \cref{sec:undirected}, the numerical facts used in the
proofs, and the final step: from a family of separated checks and a tower of quotients carrying codes of rate $\rho>1/2$
annihilated by the descended checks, the Lean library proves the negation of the eleventh paper's formal statement of
\cref{prob:q91}, \lean{MergeDepth.Q91At}~\cite{MergeDepth2026}, and, given in addition an injective encoding of the field
in $b$ bits with $\rho\log_2\abs\F/b>1/2$, the negation of its binary form \lean{MergeDepth.Q91BinaryAt}. For the
symmetrized graphs and their blow-ups it proves that, for every $K$, eventually the degrees are bounded, no closed walk of
odd length is at most $K$, and the codes are storage codes of rate at least $\rho$, respectively $\rho\log_2\abs\F/b$;
the fourth paper's Question~6.1 has no formal statement, so its negation is not formalized. Not formalized are the probabilistic step of \cref{thm:D}
(Bernstein's inequality and a union bound), the passage from a lift selection on a finite graph to separated checks on
the tree (covering theory and transport), the annihilation of quotient Tanner codes by the descended checks, the
counting of balls in the tree, the Reed--Solomon local codes, the icosian group and its congruence subgroups, the girth
bound, the whole arithmetic \cref{sec:arith}, which uses Riemann--Roch on compactified Hilbert modular surfaces in
characteristic $p$, partial Hasse invariants and the uniformization of the supersingular locus, and Theorem~A$^*$;
none of the geometry is available in Mathlib.

\paragraph{Related work.} Half bounds for storage codes are known for bipartite side information, where they follow from a
cut argument~\cite{ChristofidesMarkstrom2011,Mazumdar2015}, for odd cycles~\cite{AtkinsRombachSkerman2017}, and for
side-information graphs with a partition into two independent sides~\cite[Theorem~9]{MazumdarMcGregorVorotnikova2019}; the
fifth paper proves them for relations whose incidence components separate the two copies of every coordinate, for
hyperfinite relations, and for complex-linear codes with bounded coefficients~\cite{DirectedRecovery2026}. Rates above one
half, even tending to one, are known on graphs without short odd cycles at unbounded degree: Golovnev and Haviv construct
generalized Kneser graphs with no short odd cycles and small minrank over every finite field~\cite[Theorem~1.10]{GolovnevHaviv2022},
refuting a conjecture of Codenotti, Pudl\'ak and Resta~\cite{CodenottiPudlakResta2000}, and Haviv treats the real
case~\cite{Haviv2019}. Their degrees grow with the number of vertices, and disjoint copies of a fixed graph have bounded
degree but bounded paired distances. We found no earlier example of rate above one half at bounded degree with growing
paired distances, or on graphs of bounded degree without short odd cycles. Algebraic geometry has supplied good codes since the work of Tsfasman, Vl\u{a}du\c{t} and
Zink~\cite{TsfasmanVladutZink1982}; Tanner codes on Ramanujan graphs~\cite{LPS1988} with Reed--Solomon local codes are
classical, and Kopparty and Tamo have recently studied algebraic expander codes~\cite{KoppartyTamo2026}. We know of no
earlier use of Goren--Oort strata for codes. The twelfth paper of the series uses the same Goren--Oort Tanner operators at
$p=3$ to show that L\"uck approximation fails in positive characteristic~\cite{Luck2026}; Theorem G is the analogous
dimension bound for every inert prime $p\ge7$ and every weight $2\le k\le p-2$, and \cref{sec:arith} follows that
paper.

\paragraph{Organization.} \Cref{sec:prelim} recalls the question and sets up graphs and Tanner codes. \Cref{sec:lifts}
proves Theorem D, \cref{sec:peeling} constructs private checks and separated structures, and \cref{sec:descent} proves the
descent and Theorem T. \Cref{sec:arith} proves Theorem G. \Cref{sec:theoremA} assembles Theorems A and A$^*$, and
\cref{sec:binary} the binary Corollary A$'$; \cref{sec:undirected} answers the undirected crux (Corollary U).
\Cref{sec:limits} discusses what the counterexample changes for lower bounds
and constructions, and the limits of the method; \cref{sec:open} lists open problems. The appendices list the Lean
declarations and explain how to build and check them.
